\documentclass[10pt,twocolumn]{article}
\usepackage[T1]{fontenc}
\usepackage[utf8]{inputenc}
\usepackage{lmodern}
\usepackage[margin=1.8cm,top=2.4cm,bottom=2cm,columnsep=0.6cm]{geometry}
\usepackage{fancyhdr}
\usepackage{titlesec}
\usepackage{booktabs}
\usepackage{amsmath,amssymb,amsfonts}
\usepackage{microtype}
\usepackage{xcolor}
\usepackage{mdframed}
\usepackage{parskip}
\usepackage{enumitem}
\usepackage{hyperref}

\definecolor{rulered}{RGB}{180,30,30}
\definecolor{boxgray}{RGB}{245,245,245}
\definecolor{keywordgray}{RGB}{100,100,100}

\pagestyle{fancy}
\fancyhf{}
\fancyhead[L]{\textsc{\small Claude Mag}}
\fancyhead[C]{\small\textit{Machine Learning Research Digest}}
\fancyhead[R]{\small 2026-11-16}
\fancyfoot[C]{\small\thepage}
\renewcommand{\headrulewidth}{0.8pt}

\titleformat{\section}{\normalsize\bfseries\color{rulered}}{}{0pt}{}%
  [\vspace{-4pt}\color{rulered}\rule{\columnwidth}{0.4pt}]
\titlespacing{\section}{0pt}{8pt}{4pt}

\hypersetup{colorlinks=true,urlcolor=rulered,linkcolor=black}

\begin{document}
\setlength{\parindent}{0pt}
\setlength{\parskip}{4pt}

{\small\color{keywordgray}\textbf{Keywords:}
  Vision-Language Models \textbullet{} Multimodal Reasoning \textbullet{} KV Cache \textbullet{} Implicit Attention \textbullet{} NeurIPS 2026}\par\vspace{4pt}

{\LARGE\bfseries\raggedright Look Before You Reason: Implicit Visual Thinking for Efficient Multimodal Reasoning}\par\vspace{4pt}
{\large\itshape\raggedright Teaching VLMs to Silently Identify Decision-Relevant Image Regions Before Generating a Single Word}\par\vspace{6pt}
{\small\textbf{By} Research team (NeurIPS 2026 Poster)}\par\vspace{2pt}
{\small\color{keywordgray}arXiv:2610.27341 \textbullet{} cs.CV, cs.AI \textbullet{} NeurIPS 2026}
\par\vspace{2pt}\noindent{\color{rulered}\rule{\columnwidth}{0.6pt}}\vspace{6pt}

Before a vision-language model (VLM) generates its first reasoning token, it should already know \emph{where} in the image to focus.
The LoRe (Look-before-Reason) framework teaches VLMs to predict the most relevant image region and amplify the corresponding key-value cache entries so the decoder naturally attends to those patches throughout reasoning---no extra tool calls, no additional image inputs, no multi-turn interaction.
Accepted as a Poster at NeurIPS 2026.

\begin{mdframed}[backgroundcolor=boxgray,linecolor=rulered,linewidth=1pt,
    innerleftmargin=6pt,innerrightmargin=6pt,
    innertopmargin=6pt,innerbottommargin=6pt]
\textbf{\small AT A GLANCE}\par\vspace{4pt}
\begin{description}[leftmargin=0pt,labelsep=4pt,itemsep=2pt,parsep=0pt,
    font=\small\bfseries]
  \item[Problem:] VLM decoders must attend to relevant patches without guidance; external tools add latency; multi-turn inspection complicates training.
  \item[Why it matters:] Fine-grained visual reasoning (medical imaging, robotics, dense scenes) requires precise visual focus that standard patch encoding does not guarantee.
  \item[Method:] Lightweight region-prediction head outputs a bounding box; KV cache entries for enclosed patches are amplified; RL aligns region quality with answer accuracy.
  \item[Result:] +2.5 on MMBench, +2.4 on ScienceQA, +7.2 on COCO-ObjQA.
  \item[Evidence:] Three multimodal benchmarks; new COCO-ObjQA benchmark released; \textless{}5\% inference overhead.
\end{description}
\end{mdframed}

\section{The Big Picture}

Modern vision-language models encode an image as a grid of patch embeddings appended to the token sequence.
The decoder must attend to all patches simultaneously via self-attention, with no prior knowledge of which patches carry task-relevant information.
For simple visual questions---``What color is the car?''---this works well: the relevant region is large and prominent.
For fine-grained reasoning---``Which of the labeled containers shows a greater volume?''---the relevant region may be a small sub-image, indistinguishable from background at the patch embedding level.

The natural solution is to let the model \emph{look first}, then reason.
But existing implementations of this idea rely on external detection models, zoom operations, or OCR tools---adding significant inference overhead, complicating deployment, and introducing failure modes when tool calls go wrong.

LoRe asks: can a VLM learn to focus attention \emph{internally}, without any external machinery?

\section{Why Existing Approaches Fall Short}

\textbf{Explicit visual tools.}
Methods like ``thinking with images'' use the model to generate a sequence of tool calls (detect, crop, zoom, OCR) that progressively inspect the relevant parts of the image.
These are powerful but add latency proportional to the number of tool calls and require the model to generate syntactically valid tool invocations---a training challenge distinct from reasoning quality.

\textbf{Multi-turn visual inspection.}
Some approaches re-query the image in a second turn after generating an initial assessment.
This extends context length and requires specialised multi-turn training pipelines.
It also constrains deployment to systems that support interactive reasoning.

\textbf{Coarse attention biasing.}
Prior work has explored injecting visual saliency maps as soft position biases.
These are computed from visual features alone and do not condition on the \emph{question}, missing the fundamental point that relevance is question-dependent.

\section{Core Mechanism}

LoRe introduces two components that work together:

\begin{enumerate}[leftmargin=1.5em,itemsep=2pt]
  \item \textbf{Region Prediction Head:} A lightweight MLP on top of the visual encoder that outputs a bounding box $(x_1,y_1,x_2,y_2)$ conditioned on the question.
    It is computed \emph{before} the decoder generates any tokens.
  \item \textbf{KV Cache Amplification:} Key and value vectors for image patches whose positions fall within the predicted bounding box are scaled by a constant $\alpha > 1$.
    This soft amplification steers decoder attention toward the predicted region throughout the entire reasoning chain.
\end{enumerate}

Both components are trained end-to-end with a joint reinforcement learning objective that rewards region quality \emph{through} downstream answer accuracy.

\section{Technical Formulation}

\subsection*{Region Prediction Head}

Let $\mathbf{V} = [\mathbf{v}_1, \ldots, \mathbf{v}_P] \in \mathbb{R}^{P \times d}$ be the patch embedding matrix from the visual encoder.
Let $\mathbf{q} \in \mathbb{R}^{d_q}$ be the pooled question representation from the language model's first few transformer layers.
The region prediction head is:

\begin{equation}
\hat{b} = \mathrm{MLP}\bigl(\mathrm{pool}(\mathbf{V}),\; \mathbf{q}\bigr) \in [0,1]^4
\end{equation}

where pool is global average pooling over patch embeddings and the MLP has two hidden layers of dimension $d$.
The output $\hat{b} = (\hat{x}_1, \hat{y}_1, \hat{x}_2, \hat{y}_2)$ represents normalized image coordinates.

The supervised training loss combines L2 regression and IoU loss:

\begin{equation}
\mathcal{L}_{\mathrm{region}} = \|\hat{b} - b^*\|_2^2 + \bigl(1 - \mathrm{IoU}(\hat{b}, b^*)\bigr)
\end{equation}

where $b^*$ is the ground-truth bounding box for the question-relevant region.

\subsection*{KV Cache Amplification}

Identify the set of patches $\mathcal{P}_{\hat{b}}$ that fall within predicted box $\hat{b}$:

\begin{equation}
\mathcal{P}_{\hat{b}} = \{p : \mathrm{center}(p) \in \hat{b}\}
\end{equation}

For each patch $p \in \mathcal{P}_{\hat{b}}$ and each transformer layer $l$, amplify its key-value representations:

\begin{equation}
K^{(l)}_p \leftarrow \alpha \cdot K^{(l)}_p, \qquad V^{(l)}_p \leftarrow \alpha \cdot V^{(l)}_p
\label{eq:amplify}
\end{equation}

with amplification factor $\alpha \in [1.5, 2.0]$.
Patches outside $\mathcal{P}_{\hat{b}}$ are unchanged.
This manipulation is performed once, before any decoder tokens are generated, and persists throughout the reasoning chain.

The effect on attention weights is tractable:
given attention logits $e_{t,p} = (\mathbf{q}_t)^\top K_p / \sqrt{d_k}$, the scaled logit for $p \in \mathcal{P}_{\hat{b}}$ becomes $e_{t,p} + \log\alpha$.
Softmax normalization then increases the attention weight on amplified patches by a factor of approximately $\alpha$ (in the limit of many patches), without zeroing non-amplified patches.

\subsection*{RL Fine-Tuning}

After supervised pre-training of the region head, the full model is fine-tuned with GRPO:

\begin{equation}
\mathcal{R}(q, \hat{b}, \hat{a}) = \underbrace{\mathbf{1}[\hat{a} = a^*]}_{\text{answer reward}} + \lambda_{\mathrm{iou}} \cdot \underbrace{\mathrm{IoU}(\hat{b}, b^*)}_{\text{region reward}}
\label{eq:rl}
\end{equation}

GRPO generates a group of $G$ responses per question and normalises the reward within the group before computing the policy gradient.
The joint reward in Eq.~(\ref{eq:rl}) aligns region quality with answer correctness: the model is rewarded only when it both identifies the right region \emph{and} answers correctly.
Over training, this causes the region predictor to learn which visual context is causally relevant to the answer, not just visually salient.

\section{Learning or Inference Procedure}

\textbf{Training stages:}
\begin{enumerate}[leftmargin=1.5em,itemsep=1pt]
  \item \textbf{Visual grounding pre-train:} Train region head with $\mathcal{L}_{\mathrm{region}}$ on a visual grounding dataset (RefCOCO+). Decoder weights frozen.
  \item \textbf{Joint SFT:} Fine-tune region head and decoder jointly on question-answer pairs with grounding annotations.
  \item \textbf{RL fine-tuning:} Fine-tune with GRPO using reward Eq.~(\ref{eq:rl}) on 50K QA pairs without grounding annotations.
\end{enumerate}

Stage 3 is the most important: it aligns region prediction with answer quality using only binary answer feedback, making the method applicable to benchmarks where ground-truth bounding boxes are unavailable.

\textbf{Inference:} Standard single-pass generation.
The region head runs once ($\approx$3\% of forward-pass FLOP); KV amplification is a free element-wise multiply.
Total overhead: $<$5\% per token on average.

\section{What the Guarantee Says}

Under the cross-attention view of the VLM, Eq.~(\ref{eq:amplify}) is equivalent to adding a constant $\log\alpha$ to the attention logits of the predicted region.
This is a form of \emph{attention bias injection}, which has been shown to shift the effective attention distribution without altering the model's expressive power.
The model retains the ability to attend to non-amplified patches; amplification is a soft prior, not hard masking.

The RL training guarantee (from GRPO theory) is that the policy converges to a local optimum of the expected reward Eq.~(\ref{eq:rl}) under mild regularity conditions on the reward landscape.

\section{Experimental Findings}

\begin{center}
\small
\begin{tabular}{lrrr}
\toprule
\textbf{Method} & \textbf{MMBench} & \textbf{SciQA} & \textbf{ObjQA} \\
\midrule
Base VLM & 80.7 & 89.4 & 69.2 \\
+ Zoom tool & 81.9 & 90.3 & 72.8 \\
+ Multi-turn & 82.1 & 90.6 & 73.4 \\
LoRe (ours) & \textbf{83.2} & \textbf{91.8} & \textbf{76.4} \\
\bottomrule
\end{tabular}
\end{center}

LoRe outperforms both tool-based and multi-turn baselines on all three benchmarks, with the largest advantage on COCO-ObjQA, the benchmark most requiring precise object-level visual focus.

\section{Ablations and Interpretation}

\begin{itemize}[leftmargin=1em,itemsep=2pt]
  \item \textbf{Region only (no amplification):} +1.1 on MMBench. Region embeddings are richer, but attention is unfocused.
  \item \textbf{Amplification with random region:} $-$0.4 on MMBench. Misfocusing attention hurts.
  \item \textbf{Full LoRe:} +2.5 on MMBench.  Both components are necessary.
  \item \textbf{Amplification factor $\alpha$:} $\alpha=1.5$--$2.0$ is optimal; $\alpha > 3$ overly suppresses non-amplified patches and hurts accuracy by 0.9 on MMBench.
  \item \textbf{Region predictor quality:} IoU with ground truth correlates with final accuracy improvement, confirming that region quality drives the gain.
\end{itemize}

\section*{Reference}
(NeurIPS 2026 Poster).
\textbf{Look Before You Reason: Implicit Visual Thinking for Efficient Multimodal Reasoning.}
arXiv:2610.27341. cs.CV, cs.AI. NeurIPS 2026.
\url{https://arxiv.org/abs/2610.27341}

\end{document}
